\section{存在性与唯一性}
\label{sec:01-Mathematical-Language/03-Proof-Techniques/03}

\subsection{你的答案真的完整吗}

来，先做一道看起来极其简单的题：

\begin{center}
解方程 \(x^2=4\)。
\end{center}

你写下了 \(x=2\)，心想这还用想吗。代入验证：\(2^2=4\)，没问题。

但同桌写的是 \(x=-2\)。谁对？

都对了——可也都只说对了一半。方程 \(x^2=4\) 有两个解，只写出 \(2\) 不算完整答案，只写出 \(-2\) 也不算。完整的回答是 \(x=2\) 或 \(x=-2\)。

这不是什么深奥的数学，但这件事恰好暴露了一个极容易踩的坑：当我们说``解方程''的时候，其实同时下了两个承诺。第一个承诺是``我找到的东西确实是解''——这就是存在性。第二个承诺是``我没有漏掉任何解''——这就是唯一性。多数人只完成了第一个承诺，第二个根本就没意识到它存在。

\subsection{一个符号拆成两项独立任务}

数学里有个专门做这件事的符号：\(\exists!x\in D,P(x)\)，读成``存在唯一的 \(x\) 在 \(D\) 中满足 \(P(x)\)''。它看上去是一个符号，却在逻辑上分解成两个必须分别验证的断言：

\begin{enumerate}
\tightlist
\item
  存在性：至少有一个 \(x\in D\) 使得 \(P(x)\) 成立。你得把见证交出来，或者论证它一定藏在某个地方。
\item
  唯一性：如果 \(x_1\) 和 \(x_2\) 都满足 \(P\)，那么 \(x_1=x_2\)。任何两个候选都必须是同一个东西。
\end{enumerate}

注意唯一性的说法——它并不要求你``找到所有解然后数一下是不是只有一个''。它的逻辑形式是：随便抓两个满足条件的对象，证明它们只能相等。这个区别在无穷论域里尤其关键，因为你没法真的列出所有解再去数。

\subsection{构造性存在：见证请直接拿出来}

最直接的存在性证明就是直接把符合条件的东西写出来。

证明：任意奇数都能写成两个相邻整数之和。

任取一个奇数 \(n\)，它一定可以写成 \(n=2k+1\) 的形式，其中 \(k\) 是整数。那就直接写：

\[
n = k + (k+1).
\]

\(k\) 和 \(k+1\) 确实是相邻整数。证完了。

不光是证明了存在，而且告诉了你具体怎么找——给定任何一个奇数 37，你马上知道 \(37=18+19\)。给定 101，答案是 \(101=50+51\)。构造里每个符号都对应着可执行的步骤。

但构造完之后得逐一对照题目里的每一个限定条件。如果题目说``正整数 \(x,y\)''而你构造出了 \(x=0\) 或 \(x=-1\)，证明就有漏洞。如果题目说``两个不同的对象''而你造出来的是同一个东西用了两次，那也不行。构造性证明的最后一公里是条件逐条校对。

\subsection{非构造存在：东西肯定有，但先别管在哪儿}

也有大量定理只承诺存在，却不告诉你具体是哪个。这不是数学家在耍赖——有时候``存在''本身就是一个足够有力的结论。

连续函数的介值定理是经典的例子：设 \(f\) 在 \([0,1]\) 上连续，且 \(f(0)<0<f(1)\)，那么一定存在某个 \(c\in(0,1)\) 使得 \(f(c)=0\)。直觉上一目了然——一条不抬笔的曲线从 \(x\) 轴下方出发，终点在 \(x\) 轴上方，中间肯定得穿过 \(x\) 轴至少一次。但这个 \(c\) 到底在哪个精确位置上，定理没有公式给你。严格证明留到微积分部分，这里用它展示一种思维：知道某个东西一定存在，本身已经是一块有力的推理积木。

更漂亮的例子是下面这个经典的非构造存在论证。

\begin{center}
证明：存在无理数 \(a,b\)，使得 \(a^b\) 是有理数。
\end{center}

我们不知道具体是哪一对——但这不碍事。取 \(a=b=\sqrt{2}\)。现在分两种情况：

\begin{itemize}
\tightlist
\item
  情况 1：\((\sqrt{2})^{\sqrt{2}}\) 是有理数。那就直接找到了——\(a=b=\sqrt{2}\) 这一对就是见证。
\item
  情况 2：\((\sqrt{2})^{\sqrt{2}}\) 是无理数。那就令
  \[
  x = (\sqrt{2})^{\sqrt{2}}, \qquad y = \sqrt{2}.
  \]
  此时 \(x\) 和 \(y\) 都是无理数，而
  \[
  x^y = \bigl((\sqrt{2})^{\sqrt{2}}\bigr)^{\sqrt{2}} = (\sqrt{2})^{2} = 2,
  \]
  是一个有理数。
\end{itemize}

两种穷尽了一切可能的情况，每种都能给出见证。论证结束后，你仍然不知道 \((\sqrt{2})^{\sqrt{2}}\) 到底是不是有理数，也不影响结论的严谨。存在性被证明了，但见证的身份悬在半空——这就是非构造存在的标志性风格。这个论证同时展示了分类讨论的力量：把可能性空间切成互斥且穷尽的几块，每一块内部都有一份见证。

非构造不等于含糊。每个分支的逻辑覆盖、每个运算的合法性，都必须像构造性证明一样经得起审查。它只是提供的信息量少了一些——你知道抽屉里有东西，但不知道在哪个抽屉。

\subsection{唯一性的标准武器：假设有两个}

唯一性证明有一个几乎万能的开场白：``假设存在两个满足条件的对象 \(x_1\) 和 \(x_2\)，我们证明 \(x_1=x_2\)。''

比如证方程 \(ax=b\) 在 \(a\neq 0\) 时有唯一实解。存在性不用多说，\(x=b/a\) 代进去就成立。唯一性呢？假设 \(x_1\) 和 \(x_2\) 都满足方程：

\[
ax_1 = b = ax_2.
\]

两式相减：\(a(x_1-x_2)=0\)。因为 \(a\neq 0\)，必须 \(x_1=x_2\)。

这里没有去算两个解各自等于什么值再比较，而是直接研究了两个候选之间的差异，然后利用方程的结构（\(a\) 非零这个条件）把差异压成了零。``假设两个，推导它们相等''这个模板在不同数学领域里反复出现，只是换了一层皮：

\begin{itemize}
\tightlist
\item
  线性代数：证 \(Ax=b\) 有唯一解，等价于证 \(Ax=0\) 只有零解——两个解的差正好满足齐次方程。
\item
  微分方程：两个解的差满足对应的齐次方程，初始条件为零时这个差恒为零。
\item
  优化：严格凸函数最多有一个最小值点——如果有两个，在它们之间取中点，凸性会给出更小的函数值，矛盾。
\end{itemize}

各领域的术语不同，但底层逻辑是同一条：取出两个候选，算出差异，然后用结构的性质把差异消灭。这就是唯一性证明的核心肌肉记忆。

另一种等效的写法是：先构造出一个候选 \(x_0\)，然后证明任何满足条件的东西，推导到最后都必须等于 \(x_0\)。无论选哪种形式，量词结构不变——任取两个候选，证它们相等。

\subsection{小心：这两个方向并不互相送}

有些命题存在性容易、唯一性难；有些反过来。

\(x^2=1\) 在实数里存在性毫无问题——\(x=1\) 就是解——但唯一性不成立，因为 \(x=-1\) 也是解。存在了，但不是只有一个。

\(x^2=-1\) 在实数里没有任何解。这时候连谈唯一性的资格都没有——一旦存在性失败，\(\exists!\) 直接为假，不需要再关心唯一性那条路径。

更具欺骗性的是自然语言里悄悄夹带的唯一性预设。``最大的素数''——这个短语暗示了存在且唯一。但实际上素数无限多，根本不存在最大的那个。谁要是从``设最大的素数为 \(p\)''开始推导，后面的每一步推理都在围绕一个不存在的对象打转，写得再漂亮也是空转。

日常表述里的``这个''``那个''``唯一的''，落笔成数学之前，先停一秒确认：它真的存在吗。这一步省不得。
